% ===========================================================================
%  CAPITULO 4 — Tecnologia, IA, Segurança e Operação (T01–T11)
% ===========================================================================
\chapter{Tecnologia, IA, Segurança e Operação}
\label{cap:tecnologia}

% ---------------------------------------------------------------------------
\section{T01 · Navegador, dispositivo e conectividade}
\label{sec:t01}

\begin{resposta}{T01}{Combinações suportadas}{\parcial}
\begin{campos}
  \item[Hoje] A plataforma é 100\% web: o hub é uma aplicação Next.js e cada
    missão é um pacote estático servido por CDN, sem instalador. As suítes
    automatizadas de ponta a ponta rodam no \textbf{Chromium} (Playwright
    1.62), em 1280$\times$950 (estudante) e 1440$\times$1000 (professor). A
    responsividade do hub foi medida nos seis papéis em \textbf{375, 768 e
    1280\,px}, em cerca de 80 rotas, e as telas que rolavam de lado foram
    corrigidas (\adrr{ADR-159}). Missões desenhadas para celular declaram
    \cd{viewport} no manifesto e rodam numa moldura de 375$\times$812 dentro do
    hub. Os arquivos das missões têm caminho imutável por versão e ficam em
    cache na CDN; só a página de entrada expira em 60\,s.
  \item[Limitações] Não há matriz formal de compatibilidade: Firefox, Safari
    (WebKit) e Edge não são testados automaticamente, e não há versão mínima
    declarada nem sistema operacional homologado. Não há orçamento de tamanho
    de pacote nem teste sob banda reduzida. A medição de responsividade não
    cobriu modais nem telas com dados reais de aula, e o script dela não está
    versionado. Algumas missões foram desenhadas para computador (o Jornal da
    Escola, por exemplo, não declara \cd{viewport}). A Vila do Lume (jogo 3D)
    exige WebGL, teclado e mouse e não tem controle por toque.
  \item[Planejado] Nada registrado além do gatilho da \adrr{ADR-159} (tela nova
    que role de lado em 375\,px).
  \item[Evidência] \adrr{ADR-159}; \arq{scripts/e2e/lib/roteiro.mjs};
    \arq{scripts/\B{}e2e/\B{}estudante/\B{}16-assistente-invisivel.mjs} (emulação de
    celular 412$\times$915); RNF-03 em \cd{PU-DOC-001}, Cap. 3.
  \item[Lacuna] Engenharia precisa executar e anexar a matriz (Chrome, Edge,
    Firefox e Safari recentes; Android e iOS; Chromebook) antes de o TR fixar
    versões mínimas.
\end{campos}
\end{resposta}

\begin{paraotr}
A solução deve funcionar, sem instalação de software, nas duas últimas versões
estáveis de Chrome, Edge, Firefox e Safari, em telas a partir de 360\,px de
largura para as funções do estudante e a partir de 1280\,px para a condução da
aula. Aceite: roteiro de homologação executado pela contratada em cada
combinação, com registro de tela e lista de restrições por missão.
\end{paraotr}

% ---------------------------------------------------------------------------
\section{T02 · Autenticação e identidade da rede}
\label{sec:t02}

\begin{resposta}{T02}{Login, vínculo e revogação}{\parcial}
\begin{campos}
  \item[Hoje] O login em produção é por \textbf{usuário e senha próprios} da
    plataforma: senha guardada com \emph{hash} bcrypt, limite de 8 tentativas
    por minuto por usuário, mensagem de recusa genérica (\adrr{ADR-97}) e troca
    obrigatória da senha provisória. Depois do login, a plataforma emite
    \emph{tokens} assinados (RS256, validade de 15 minutos) com usuário,
    escola, turma principal e papel, e publica a chave pública num endereço
    JWKS padrão; a API valida emissor, audiência e assinatura. Esse desenho
    segue o padrão \textbf{OpenID Connect}, de modo que trocar o emissor por um
    provedor de identidade da rede é troca de configuração. Escola, turma e
    papel vêm das lotações e matrículas cadastradas pela importação (T03);
    quem trabalha em mais de uma escola escolhe a escola ativa no cabeçalho,
    validada contra as lotações. A sessão expira após 2 horas sem uso.
  \item[Limitações] A federação com o provedor da rede (WSO2 Identity Server,
    via OIDC) está \textbf{preparada mas não operacional}: o conector existe e
    liga por configuração, mas o mapeamento dos atributos da rede ainda não
    foi feito. \textbf{SAML não existe atualmente.} Desativar uma conta
    impede novos logins na hora, mas uma sessão já aberta continua válida até
    expirar (até 2 horas de inatividade); não há revogação ativa nem
    desconexão global. Perder uma turma vale na hora (cada pedido confere a
    matrícula); perder a escola ou mudar de papel vale no próximo login.
  \item[Planejado] Ligar o WSO2 por contrato (\arq{docs/INFRA-IDENTITY.md}:
    um \emph{tenant} por contrato, emissor e JWKS próprios). Decisões
    abertas: base de usuários mestre ou federação com criação na primeira
    entrada (Google Workspace, Gov.br), formato do usuário e tela de login.
  \item[Evidência] \cd{PU-DOC-001}, Cap. 5, diagramas \emph{Três relógios:
    token, renovação e sessão} e \emph{A jornada do estudante, do login à
    missão}; \arq{apps/\B{}api/\B{}tests/\B{}senha-provisoria.spec.ts};
    \arq{scripts/\B{}e2e/\B{}publico/\B{}02-login-e-papeis.mjs}; \adrr{ADR-95},
    \adrr{ADR-96}.
  \item[Lacuna] Decidir se encerrar lotação ou desativar conta deve derrubar a
    sessão aberta (hoje não derruba).
\end{campos}
\end{resposta}

\begin{paraotr}
A solução deve oferecer autenticação própria e federação com o provedor de
identidade da contratante por OpenID Connect; SAML 2.0, quando exigido,
deve ser declarado como item opcional. A desativação de um usuário deve
impedir o acesso a dados em até \textit{[prazo a definir]} minutos, inclusive
em sessão já aberta. Aceite: demonstração do login federado com usuário de
teste da rede e da revogação dentro do prazo.
\end{paraotr}

% ---------------------------------------------------------------------------
\section{T03 · Cadastro de escolas, turmas e pessoas}
\label{sec:t03}

\begin{resposta}{T03}{OneRoster e importação}{\implementado}
\begin{campos}
  \item[Hoje] \textbf{OneRoster 1.2, formato CSV}: os sete arquivos do padrão
    (\cd{manifest}, \cd{orgs}, \cd{academicSessions}, \cd{classes},
    \cd{users}, \cd{enrollments}, \cd{roles}), um por envio, em modo
    \cd{bulk} ou \cd{delta}, pela tela do gestor ou por chamada autenticada à
    API. Aceita a extensão brasileira de CPF. Há também uma planilha
    simplificada (matrícula, nome, usuário, e-mail, papel, turma, série,
    turno; até 50 mil linhas). Modelos e exemplos de cada arquivo são
    baixáveis pela própria plataforma. Regras de carga: reimportar o mesmo
    lote não muda nada (idempotente, pela chave da rede e pelo CPF); troca de
    turma encerra a matrícula anterior sem apagá-la; CPF já existente em outra
    escola vira transferência (estudante) ou lotação adicional (professor);
    \cd{tobedeleted} desativa a pessoa sem apagar; uma carga completa
    \textbf{não} desativa quem ficou de fora, só os lista como candidatos;
    CPF repetido ou inválido é erro de linha; linhas válidas entram e as
    inválidas voltam com o motivo; um arquivo fora de ordem informa qual
    arquivo falta. O retorno traz contadores de criados, atualizados,
    transferidos, desativados e erros.
  \item[Limitações] \textbf{A plataforma não busca mudanças}: é a Secretaria
    (ou o gestor) quem envia as atualizações. Não existe a API REST do
    OneRoster nem sincronização agendada, e o pacote zip do padrão não é
    aceito (um arquivo por envio). Limites: 20 envios por minuto (OneRoster)
    e 5 por minuto (planilha). A senha do arquivo é ignorada.
  \item[Planejado] Nada registrado.
  \item[Evidência] \arq{apps/\B{}api/\B{}tests/\B{}oneroster-import.spec.ts} (cerca de 25
    casos, entre eles ``carga completa incompleta não desativa ninguém'' e
    ``reaplicar o mesmo lote não muda nada''),
    \arq{rostering-import.spec.ts}, \arq{rede-cpf-transferencia.spec.ts};
    suíte \arq{scripts/\B{}e2e/\B{}gestor\_contrato/\B{}02-oneroster.mjs};
    \arq{apps/\B{}api/\B{}src/\B{}rostering/\B{}AGENTS.md}.
\end{campos}
\end{resposta}

\begin{paraotr}
A contratada deve aceitar a carga de escolas, turmas, pessoas e vínculos no
padrão OneRoster 1.2 (CSV), por interface e por API autenticada, com
importação parcial, relatório de erros por linha e carga idempotente. A
atualização é de iniciativa da contratante, salvo integração adicional
pactuada. Aceite: carga de teste sintética com inclusão, troca de turma,
transferência, desligamento e linhas inválidas, conferida contra o relatório.
\end{paraotr}

% ---------------------------------------------------------------------------
\section{T04 · Cada usuário vê só o que lhe compete}
\label{sec:t04}

\begin{resposta}{T04}{Perfis e segregação de dados}{\implementado}
\begin{campos}
  \item[Hoje] Seis perfis: \textbf{estudante} (a própria matrícula, nunca
    gabarito nem dados de colega); \textbf{professor} (só as turmas em que
    leciona); \textbf{gestor da escola} (a própria escola);
    \textbf{gestor regional} (as escolas da regional, uma por vez, e o
    panorama somado); \textbf{gestor do contrato} (a rede inteira);
    \textbf{administração Unooria} (o anterior, mais disponibilidade de
    recursos, quotas e custos de IA). Ninguém cria o próprio nível nem um
    acima (\adrr{ADR-80}). A autorização é aplicada \textbf{na API, em toda
    chamada}, em três travas: o papel exigido pela rota; o isolamento por
    escola no próprio banco (\emph{Row-Level Security}), com o escopo tirado
    do \emph{token} assinado e não do parâmetro; e, para o professor, a
    conferência da matrícula na turma a cada pedido. Um identificador de
    outra escola manipulado na chamada cai no isolamento do banco ou volta
    403. Exportações e painéis usam as mesmas rotas protegidas. Cada contrato
    tem a própria infraestrutura (isolamento físico entre redes). Quem
    trabalha em várias escolas acumula lotações e troca a escola ativa;
    transferência de estudante encerra lotação e matrícula de origem, e o
    professor de origem deixa de vê-lo.
  \item[Limitações] Três tabelas de cadastro (pessoas, turmas e lotações)
    ficam fora do isolamento do banco: nelas, a trava é só o filtro por escola
    na consulta, e uma falha dessas já foi encontrada e corrigida
    (\adrr{ADR-49}). A revogação de sessão tem o atraso descrito em T02.
  \item[Evidência] Matriz papel $\times$ rota $\times$ tela em \cd{PU-DOC-001},
    Cap. 7 (``Nível 3''); testes \arq{tenancy.spec.ts} (``escola A não alcança
    linha da escola B nem por chave direta''), \arq{matriculas-tenancy.spec.ts},
    \arq{professor-por-matricula-tenancy.spec.ts}, \arq{planos-tenancy.spec.ts},
    \arq{rede-cpf-transferencia.spec.ts}, além de 32 especificações que
    provam respostas 403; suítes e2e de gestor de escola (quem trabalha em
    duas escolas) e de contrato (transferência e perfil portátil).
\end{campos}
\end{resposta}

\begin{paraotr}
A solução deve aplicar a autorização no servidor, em toda requisição,
impedindo que um usuário acesse dados de escola, turma ou estudante fora do
seu vínculo, inclusive por manipulação de identificadores. Aceite: roteiro de
teste de segregação executado pela contratada com dois usuários de escolas
diferentes, incluindo chamadas diretas à API com identificadores alheios, e
relatório de teste de intrusão independente antes da produção.
\end{paraotr}

% ---------------------------------------------------------------------------
\section{T05 · Conteúdo inadequado nas duas modalidades de IA}
\label{sec:t05}

\begin{resposta}{T05}{Guardrails}{\parcial}
\begin{campos}
  \item[Hoje] Toda chamada de IA passa por uma camada única no servidor;
    nenhuma instrução de sistema chega ao navegador. \textbf{Geração que o
    professor revisa}: toda instrução de geração carrega um bloco de regras
    invioláveis (pessoas fictícias; sem sexo, violência gráfica, ódio,
    drogas, automutilação ou palavrões; sem estudante ou escola reais); o
    resultado passa por validação automática com listas de termos e perfis
    proibidos; o texto do professor entra como dado delimitado e com tamanho
    máximo, uma barreira contra instruções injetadas; e o professor aprova
    antes de a turma ver (P04). \textbf{Estudante falando com a IA}: nas
    missões em que há texto livre, a IA \emph{classifica} a mensagem numa
    lista fechada de intenções e a fala devolvida vem de um catálogo escrito
    pela equipe; com confiança abaixo de 0,7, lentidão, falta de saldo ou
    recusa, vale uma árvore de regras determinística. Pedidos como ``ignore as
    instruções'' são detectados e tratados como fora do tema sem chegar à IA.
    A série é tirada da turma do estudante, nunca do pedido. Há quota diária
    por perfil e limite de chamadas. \textbf{Temas sensíveis}: o check-in de
    bem-estar é o caso mais cercado (itens fixos, bloqueio de vocabulário
    clínico e de links, sinal de risco pula a IA e mostra texto fixo com CVV
    188 e SAMU 192). \textbf{Dados pessoais}: na missão Manifesto, CPF,
    telefone, e-mail, CEP, cartão e @perfil são detectados e o texto é
    recusado sem ser gravado. Golpes, \emph{fake news} e \emph{cyberbullying}
    são tratados como objeto de estudo dentro do roteiro da missão.
  \item[Limitações] \textbf{Não há configuração própria dos filtros de
    segurança do provedor}: valem os padrões do Gemini, e a recusa por
    segurança é só contada nas métricas, sem tratamento específico. Não há
    moderação dedicada de entrada e saída. O tutor genérico devolve o texto do
    modelo sem validação de saída. O filtro de dados pessoais existe só numa
    missão e não detecta nomes próprios. \textbf{Não há registro formal de
    incidentes de IA.}
  \item[Provedor $\times$ Unooria] Do provedor: os filtros padrão do modelo. Da
    Unooria: instruções de sistema, regras invioláveis, classificação fechada,
    validação da saída, detecção de dados pessoais, quotas e a aprovação
    humana.
  \item[Evidência] \arq{apps/\B{}api/\B{}src/\B{}geracao/\B{}prompts/\B{}guardrails.ts};
    \arq{apps/\B{}api/\B{}src/\B{}manifesto-ia/\B{}regras/\B{}moderacao.ts};
    \adrr{ADR-191}--\adrr{ADR-193}; \cd{PU-DOC-001}, Cap. 11 (guardrails do
    check-in); testes \arq{ai-tutor-prompt.spec.ts},
    \arq{manifesto-ia-regras.spec.ts}, \arq{bemestar-regras.spec.ts}; conjunto
    de referência do classificador da 2ª série (88/89 e 89/89 acertos, zero
    crédito indevido).
  \item[Lacuna] Configurar e versionar os filtros do provedor; estender a
    detecção de dados pessoais a toda entrada livre de estudante; criar
    registro e rito de incidentes.
\end{campos}
\end{resposta}

% ---------------------------------------------------------------------------
\section{T06 · Dados pessoais e fluxos para a IA}
\label{sec:t06}

\begin{resposta}{T06}{Mapa de dados}{\parcial}
\begin{campos}
  \item[Hoje] A rede é \textbf{controladora} e a Unooria \textbf{operadora};
    a base legal é a execução de política pública, não o consentimento
    (página \rota{/privacidade}, \adrr{ADR-70}). Categorias tratadas:
    cadastro (nome, usuário, e-mail, CPF, matrícula, escola, turma),
    atividade nas missões, evidências por habilidade, textos escritos nas
    missões e, em missões específicas, conversa com a IA. Infraestrutura:
    Google Cloud em \textbf{\cd{us-east1}} (EUA), com banco gerenciado de alta
    disponibilidade e recuperação a um ponto no tempo; IA no Vertex AI
    (endpoint global) e voz ao vivo da Ajuda em \cd{us-central1}; avatar da
    Ajuda (LiveAvatar/HeyGen, EUA) só para profissionais. A transferência
    internacional é declarada na página de privacidade com as cláusulas-padrão
    da Resolução CD/ANPD 19/2024. \textbf{Para a IA vão textos sem
    identificação}: a série e o conteúdo, nunca nome, escola ou identificador
    do estudante. Os registros técnicos mascaram CPF, e-mail, nome e textos
    livres; a telemetria usa identificador opaco.
  \item[Limitações] \textbf{Não há prazo nem rotina de retenção}, de exclusão
    ou de anonimização: desligar alguém o desativa, sem apagar (questão aberta
    no Cap. 7 da \cd{PU-DOC-001}, ``decisão de produto e de contrato''). O CPF é
    guardado sem cifra adicional (\adrr{ADR-33}). O avatar da Ajuda não está
    listado entre os suboperadores da página de privacidade. A afirmação de
    que o provedor não treina modelos com os dados não cita a cláusula
    contratual correspondente. Só há diagrama de fluxo de dados do
    check-in de bem-estar; não há um diagrama geral.
  \item[Evidência] \arq{apps/\B{}hub/\B{}app/\B{}privacidade/\B{}page.tsx};
    \arq{apps/\B{}api/\B{}src/\B{}observabilidade/\B{}redaction.ts}; \arq{docs/INFRA.md};
    \cd{PU-DOC-001}, Cap. 7 (região e transferência) e Cap. 11.
  \item[Lacuna] Jurídico/DPO: definir retenção e exclusão; incluir o avatar
    entre os suboperadores; anexar os termos do Google Cloud sobre uso de
    dados (o serviço Vertex AI, por contrato, não usa dados do cliente para
    treinar modelos, mas a cláusula deve ser citada, não presumida).
    Engenharia: desenhar o diagrama geral de fluxo de dados.
\end{campos}
\end{resposta}

\begin{atencao}
\textbf{Ponto sensível para o TR}: os dados ficam fora do Brasil
(\cd{us-east1}) e a IA roda em endpoint global. Se a contratante exigir
residência no país, isso é uma mudança de infraestrutura, e o modelo de IA
usado hoje não está disponível na região de São Paulo. O TR deve dizer se
admite transferência internacional com as salvaguardas da ANPD ou se exige
residência, porque a resposta muda custo e prazo.
\end{atencao}

% ---------------------------------------------------------------------------
\section{T07 · Trilha de auditoria}
\label{sec:t07}

\begin{resposta}{T07}{Eventos auditáveis}{\parcial}
\begin{campos}
  \item[Hoje] Três mecanismos. (1) \textbf{Auditoria de recursos}: ligar ou
    desligar produto e missão, ativar ou fixar versão, alterar quota e
    configuração de IA. Cada registro guarda recurso, escola, ação, versão,
    motivo, valor anterior e novo, autor (com o nome congelado) e data/hora;
    alteração é recusada pelo banco. (2) \textbf{Histórico da evidência
    pedagógica}: cada marca, recálculo, confirmação e ajuste do professor
    gera um evento com valor anterior, valor novo, autor e data; o banco só
    aceita inclusão (alteração recusada, e o serviço não tem permissão de
    apagar). (3) \textbf{Registros técnicos} estruturados com código estável,
    identificador de correlação da requisição, escola e perfil, retidos por 30
    dias. Algumas missões guardam ainda o histórico de comandos do professor
    (quem pausou ou liberou e quando).
  \item[Limitações] \textbf{Não há auditoria de}: login e acesso (só registro
    técnico), criação, edição ou publicação de missão, aprovação de conteúdo,
    ações administrativas sobre usuários e leitura de dados sensíveis. A
    auditoria de recursos permite exclusão no banco. Nenhuma tela ou rota
    consulta o histórico da evidência; ele só é lido direto no banco.
  \item[Planejado] Destino de registros de auditoria por contrato com
    \textbf{400 dias} de retenção, reunindo eventos de autenticação,
    importação, publicação e atividade administrativa da nuvem
    (\adrr{ADR-57}/\adrr{ADR-58}, risco R-05, meta antes de nov./2026). Hoje
    está pendente.
  \item[Evidência] \arq{apps/\B{}api/\B{}src/\B{}recursos/\B{}auditoria.ts}; tabelas
    \cd{recursos\_auditoria} e \cd{evidencias\_eventos} (\cd{PU-DOC-001},
    Cap. 8); \cd{PU-DOC-001}, Cap. 17.
\end{campos}
\end{resposta}

\begin{paraotr}
A solução deve registrar, de forma não alterável pelo usuário e consultável
pela contratante, no mínimo: autenticação, alterações de cadastro e de
vínculo, publicação de atividade, aprovação de conteúdo gerado por IA e
revisão de avaliação, com ator, data/hora, objeto, valor anterior, valor novo,
justificativa (quando houver) e identificador de correlação, retidos por pelo
menos \textit{[prazo a definir]}. Aceite: exportação da trilha de um roteiro
de homologação executado com dados sintéticos.
\end{paraotr}

% ---------------------------------------------------------------------------
\section{T08 · Falha de conexão ou de dispositivo}
\label{sec:t08}

\begin{resposta}{T08}{Progresso, retomada e perda}{\parcial}
\begin{campos}
  \item[Hoje] \textbf{O estado da missão vive no servidor.} Cada jogada que
    vale é uma chamada à API, gravada na hora na sessão do estudante; não se
    espera o fim da missão para salvar. Recarregar a página, fechar o
    navegador ou trocar de computador retoma no passo salvo, comportamento
    registrado nas missões do catálogo (\cd{PU-DOC-001}, Cap. 12). A conclusão
    é gravada uma única vez por estudante e missão: repetir não cria registro
    novo. As rotas do estudante são idempotentes (repetir uma jogada devolve o
    estado atual), e algumas missões mantêm uma fila de reenvio para quedas
    curtas de rede. Os eventos de telemetria têm identificador gerado no
    navegador e são gravados sem duplicar; ficam num \emph{buffer} de até 500
    eventos, enviados em lotes a cada 10 segundos, com novas tentativas
    espaçadas até 60 segundos e envio de emergência ao fechar a aba.
  \item[Limitações] \textbf{Exige conexão contínua}: não há modo
    \emph{offline}, por desenho, porque a correção é feita no servidor e o
    gabarito nunca vai ao navegador. Uma jogada feita sem rede não é aceita
    até reconectar; o reenvio automático só existe nas missões que o
    implementaram. A fila de telemetria fica em memória e se perde se a aba
    for fechada à força ou faltar energia antes do envio; isso afeta só
    telemetria, não o progresso nem as evidências, que já foram gravados na
    jogada.
  \item[Planejado] Guardar a fila de telemetria no armazenamento do navegador
    (IndexedDB), sem prazo.
  \item[Evidência] \arq{packages/\B{}sdk/\B{}src/\B{}telemetry.ts} e
    \arq{telemetry.test.ts}; \arq{apps/\B{}api/\B{}src/\B{}progress/\B{}progresso.controller.ts};
    \cd{PU-DOC-001}, Cap. 10 (onde pode haver perda) e Cap. 12 (retomada por
    missão).
\end{campos}
\end{resposta}

\begin{paraotr}
A solução deve preservar o progresso e as respostas já enviadas pelo
estudante em caso de queda de conexão, fechamento do navegador ou troca de
equipamento, permitindo retomar do último passo registrado, sem duplicar
participação nem evidência. Aceite: roteiro em homologação com interrupção de
rede e reabertura em outro equipamento, por missão.
\end{paraotr}

% ---------------------------------------------------------------------------
\section{T09 · Escala, disponibilidade, monitoramento e suporte}
\label{sec:t09}

\begin{resposta}{T09}{O que pode ser comprovado}{\parcial}
\begin{campos}
  \item[Hoje] \textbf{Monitoramento}: registros estruturados, rastreamento
    distribuído, métricas e erros de navegador estão ativos no ambiente de
    desenvolvimento (\cd{PU-DOC-001}, Cap. 17). Lá estão no ar 4 sondas de
    disponibilidade a partir de 3 regiões, 12 alertas, 2 objetivos de nível de
    serviço (API e hub, 99,5\% sem erro de servidor em 30 dias), fila de
    mensagens mortas da telemetria e painel de operação. Objetivos definidos
    para produção: disponibilidade de hub e API de 99,5\%; 99\% das
    requisições abaixo de 800\,ms (p95); 98\% das chamadas de IA concluídas;
    99\% da telemetria ingerida em menos de 5 minutos. \textbf{Contingência}:
    banco com alta disponibilidade e recuperação a um ponto no tempo; hub e
    API voltam à revisão anterior em segundos; uma missão volta à versão
    anterior sem novo \emph{build}; versão nova pode ser testada numa escola
    antes da rede (Mapa de Produto, H10.3.1 e H10.3.2).
  \item[Limitações] \textbf{Nenhum teste de carga foi executado}. O
    \cd{PU-EST-001} é um plano (versão 0.1, 23/09/2026) com sete testes,
    critérios de aprovação e roteiro, e o risco R-01 o classifica como ``o
    maior risco técnico aberto''. A arquitetura foi dimensionada para cerca de
    10 mil estudantes e 2 a 4 mil simultâneos por contrato. Homologação e
    produção ainda não têm os alertas aplicados. \textbf{Não existe
    compromisso de suporte} (canais, horários, severidades, prazos de resposta
    e de restabelecimento) nem janela de manutenção formal. RPO, RTO e
    retenção de cópias de segurança não estão documentados.
  \item[Planejado] Executar os testes do \cd{PU-EST-001}; aplicar alertas em
    homologação e produção; alerta por consumo do orçamento de erro.
  \item[Evidência] \cd{PU-EST-001}; \cd{PU-DOC-001}, Cap. 6 (reversões), 17
    (SLOs e alertas) e 19 (R-01).
  \item[Lacuna] Comercial/Operação: definir o modelo de suporte. Engenharia:
    executar o teste de carga e declarar RPO/RTO.
\end{campos}
\end{resposta}

\begin{paraotr}
Disponibilidade mensal mínima de 99,5\% para as funções de estudante e
professor, medida por sonda externa e relatada mensalmente; tempo de resposta
p95 abaixo de 800\,ms em carga equivalente ao pico declarado pela contratante
(\textit{[número]} estudantes simultâneos); RPO e RTO declarados pela
contratada e comprovados em exercício de recuperação. Aceite: relatório de
teste de carga com metodologia e resultados antes da produção, e relatório
mensal de disponibilidade e incidentes.
\end{paraotr}

% ---------------------------------------------------------------------------
\section{T10 · Indicadores de adoção}
\label{sec:t10}

\begin{resposta}{T10}{Eventos e painéis de uso}{\parcial}
\begin{campos}
  \item[Hoje] \textbf{Relatório da diretoria} (Mapa de Produto, H07.2.1 e
    H07.2.2, \adrr{ADR-82}): por missão e por escola, ``turmas abertas'',
    ``turmas encerradas'' e ``estudantes com registro'', com exportação em
    planilha; o gestor regional vê a própria regional e o gestor do contrato
    escolhe a regional. \textbf{Cadastro de usuários}: filtro ``ativos /
    inativos / sem primeiro acesso''. \textbf{Análise de produto}: sessões,
    envolvimento e tempo de uso por atividade, só para uma missão
    (Unoovista). \textbf{Uso de IA}: custo e consumo por produto, modelo e
    escola, em vários períodos, mas só para a administração Unooria. Os
    números são lidos na hora da consulta. Eventos coletados: liberação e
    encerramento da missão na turma, sessões e jogadas por estudante,
    conclusão, telemetria das missões, primeiro acesso.
  \item[Limitações] O relatório filtra só por regional: \textbf{não há filtro
    por período, bimestre, ano/série ou turma}. \textbf{Não existe registro de
    último acesso}, nem contagem de professores que acessaram, de missões
    configuradas, iniciadas e concluídas por período, de participação ao
    longo do tempo, nem lista de pendências ou de baixa participação. Os
    períodos letivos existem no cadastro (importados do OneRoster), mas não
    são usados nos relatórios. Cadastrados e ativos se distinguem só no
    filtro de usuários, não num indicador.
  \item[Tratamento] Duplicidade: um registro de progresso por estudante e
    missão, e reimportação idempotente. Transferência: a identidade acompanha
    o estudante, o que ele produziu fica na escola de origem, e a escola nova
    vê o histórico.
  \item[Planejado] H07.3.2 ``Ver os números do dia'' (a fazer): turmas com
    missão aberta, estudantes que entraram e missões concluídas no dia, por
    regional. É o compromisso V1-01 (Capítulo~\ref{cap:v1}).
  \item[Evidência] \arq{apps/\B{}hub/\B{}app/\B{}(plataforma)/\B{}admin/\B{}relatorios/\B{}};
    \arq{apps/\B{}api/\B{}src/\B{}rostering/\B{}escolas.controller.ts}; Mapa de Produto,
    épicos E07 e E03.
\end{campos}
\end{resposta}

% ---------------------------------------------------------------------------
\section{T11 · Pré-visualização integral pelo professor}
\label{sec:t11}

\begin{resposta}{T11}{Preview da missão contextualizada}{\naoexiste}
\begin{campos}
  \item[Hoje] \textbf{Não existe atualmente.} A missão só abre para o perfil
    estudante: professor e gestor são levados de volta à própria área. O
    professor conta com a ficha da missão, o guia da aula e a revisão item a
    item do conteúdo gerado (P04), além de ``zerar a turma'' para reiniciar
    uma rodada. Algumas telas de preparo mostram prévia de textos e
    legendas, mas nenhuma permite jogar a missão.
  \item[Limitações] Como o recurso não existe, não há garantia a descrever
    sobre participação ou evidência geradas no preview.
  \item[Planejado] Não há história no Mapa de Produto nem item no
    \emph{backlog}. É o compromisso V1-02 (Capítulo~\ref{cap:v1}), onde está a
    proposta de implementação.
  \item[Evidência] \arq{apps/\B{}hub/\B{}app/\B{}exp/\B{}[expId]/\B{}page.tsx} (restrição ao
    estudante).
\end{campos}
\end{resposta}
